\section{Question: a second, independent route to $\theta=2/9$}
The Koide angle $\theta=2/9$ is empirically confirmed to seven digits and understood in FFGFT as
the rational address of a value fixed through a reference point (Dok.~292), not as a pole-mass
exactness claim (P42). Dok.~290/291 established the \emph{channel} in which $\theta$ lives: not
the symmetric magnitude functional (where $2/9$ is counter-blind), but the phase channel, with
$\chi=\pi/2$ as the skew-adjoint selector and $\delta$ as the still-open second harmonic. The
\emph{value} of $\theta$ remained a found, not a derived, number.

This document supplies a second, entirely independent route. It comes neither from the masses nor
from the Koide parametrisation, but purely from the \emph{geometry} of the doubling $6=3\oplus3'$
(Dok.~285): the embedding of FFGFT's own $C_3$ structure into the icosahedral ambient group
$A_5$. This embedding produces $\theta=2/9$ parameter-free, forward and exactly -- with no $2/9$
as input. That a second, methodologically foreign route hits the \emph{same} rational value is
the actual statement: not a translation of one into the other, but a convergence in the sense of
the established FFGFT logic, in which $\alpha$, the lepton masses and Koide act as
\emph{verifications} of distinct routes pointing at the same constants. The computations are in
\texttt{ikosaeder\_theta\_delta.py} (numpy-only, seed 20780458).

\section{The $C_3$-in-$A_5$ embedding and the full redistribution}
The three lepton modes are the Fourier eigenmodes of the $Z_3$ circulant on the $\mathbb{C}^3$
fibre (Dok.~282/288): the trivial mode $v_0=\tfrac{1}{\sqrt3}(1,1,1)$ and the two non-trivial,
Galois-conjugate modes $v_{1,2}=\tfrac{1}{\sqrt3}(1,\omega^{\pm1},\omega^{\pm2})$ with
$\omega=e^{2\pi i/3}$. The electron is a non-trivial mode; which of the two is decided solely by
the sign of the angle, not by the value computed here.

The doubling places $C_3$ inside the icosahedral group $A_5$; the five-fold generator is a
rotation $R_5$ by $72^\circ$ about an icosahedral $5$-fold axis (standard embedding: axis
$(0,1,\varphi)$ with $\varphi=\tfrac{1+\sqrt5}{2}$). Since the $3$-fold and $5$-fold axes point
in different directions, $R_5$ mixes the $C_3$ modes. The \emph{full} redistribution of the
electron mode $v_e=v_1$ has weights $p_j=\lvert\langle v_j\vert R_5\,v_e\rangle\rvert^2$:
\[
\underbrace{p_0=\frac{2}{9}}_{\text{trivial }v_0},\qquad
\underbrace{p_2=\frac{5-3\varphi}{9}}_{\text{2nd harm.\ }v_2},\qquad
\underbrace{p_1=\frac{2+3\varphi}{9}}_{\text{stays }v_e},\qquad p_0+p_1+p_2=1.
\]
All three weights carry the denominator $9=3^2$ -- the same $Z_3$ structure in which $2/9$ sits
as an orbit address. Confirmed numerically (40 digits, mpmath): $p_0=0.2222\ldots=2/9$ exactly,
difference $<10^{-30}$.

\section{The decisive finding: $\theta=2/9$ as the trivial share}
The share that, under the icosahedral $5$-fold, passes into the \emph{trivial} mode $v_0$ is
\[
\boxed{\;p_0=\big\lvert\langle v_0\,\vert\,R_5\,v_e\rangle\big\rvert^2=\frac{2}{9}\;}
\]
exactly and parameter-free. No mass ratio, no Koide parametrisation and no $2/9$ enters the
value; it follows from the $C_3$ modes and the icosahedral rotation alone. The associated
\emph{total} leak is $p_0+p_2=(7-3\varphi)/9=0.23843$, the second, parameter-free candidate for
the amplitude $\delta$ of the second harmonic; it agrees with the $\delta^\ast=0.23887$
back-computed from $\theta=2/9$ to $0.18\,\%$ (just at the $\tau$-limited tolerance of
$0.135\,\%$, $\sim1.3\sigma$). Value \emph{and} angle thus stem from the \emph{same} $5$-fold
redistribution: $\theta$ is the trivial, $\delta$ the total share. The earlier separate treatment
as single harmonics (first for $\theta$, second for $\delta$) was the too-narrow view; the full
redistribution carries both.

\section{Robustness: the specificity of the icosahedral geometry}
The value $2/9$ is not a generic property of any rotation, but of the icosahedral one:
\begin{itemize}
\item \textbf{Axis specificity.} $200$ random rotation axes at $72^\circ$ give the $v_0$ leak in
the range $[0.036;\,0.452]$; \emph{not one} hits $2/9$. Only the icosahedral $5$-fold axes do.
\item \textbf{Angle and order specificity.} The angle scan shows $2/9$ exactly at $72^\circ$; the
$n$-fold series yields $2/9$ only at $n=5$ ($2/5$ at $n=3$, $4/9$ at $144^\circ$). The rational
cleanliness occurs at the distinguished angles.
\item \textbf{Mode independence.} Both non-trivial modes $v_1,v_2$ give $p_0=2/9$; the trivial
$v_0$ would give $4/9$. The result does not depend on which conjugate mode is read as the electron.
\item \textbf{Axis choice (discrete, P3).} Of the six $5$-fold axes, three ($+\varphi$ branch,
$3$-$5$ axis angle $37.38^\circ$) give $2/9$, three ($-\varphi$ branch, $79.19^\circ$) give
$4/9$. The physical embedding $C_3<A_5$ selects the nearest $5$-fold, the $+\varphi$ branch --
hence $2/9$. It is a discrete two-way choice, not a free parameter.
\end{itemize}

\section{Two witnesses, one value}
The $\theta$ question thus has two mutually independent sources that coincide exactly:
\begin{center}
\begin{tabular}{@{}p{0.31\linewidth}p{0.31\linewidth}p{0.30\linewidth}@{}}
\toprule
\textbf{Route} & \textbf{Origin} & \textbf{Value}\\
\midrule
Koide (phenomenological) & from the measured lepton masses; no geometry & $\theta=2/9$\\
$C_3$-in-$A_5$ (geometric) & from the $5$-fold redistribution of the modes; no masses, no $2/9$ input & $p_0=2/9$\\
\bottomrule
\end{tabular}
\end{center}
The two routes know nothing of each other -- one knows only masses, the other only geometry --
and yet hit \emph{exactly} the same rational number. This is a convergence, not a translation
deficit: it is not required that the geometric route translate back into the Koide cosine
parametrisation. Just as a derived value need not be identical to its measurement prescription to
confirm it, the parameter-free geometric prediction $2/9$ need not take the shape of the
phenomenological fit parameter to agree with it. That the geometric value does \emph{not} vary
with $\theta$ is here the hallmark of a prediction, not a defect: the geometry has no free
parameter and names a fixed value.

\section{$\theta=2/9$ as a translation invariant of two models}
The \enquote{second witness} must be stated more precisely so as not to be misread as an
ontological claim. It is \emph{not} about a geometric truth beating a phenomenological one, nor
about the torus \enquote{really} being the Hilbert space or vice versa. Both -- the
point/differential geometry of the $T^4$ and the spectral geometry of the Hilbert space
$L^2(T^4)\otimes\mathbb{C}^3$ -- are \emph{mathematical models} that are demonstrably and
losslessly translatable into one another (Dok.~230/282): the $Z_3$ fibres, the $C_3$ Fourier
modes, the $A_5$ embedding are the same structure in both languages, merely represented
differently.

The actual finding is therefore methodological, not ontological: $\theta=2/9$ is an \emph{invariant
of the translation}. It appears in the spectral representation (as the redistribution weight
$p_0$) just as in the phenomenological Koide address, because it belongs to the shared structure
and not to a single representation. This is precisely what makes it robust: the constant falls out
of the structure irrespective of which of the translatable languages one computes in. This lies on
the established line in which $\hbar$ and $c$ are conversion factors and not ontological
primitives: what counts is what is invariant under the conversion, not the choice of one side as
\enquote{true}. On the question of which of the two models is reality, nothing is asserted here
deliberately -- it is not needed, and the computation does not supply it. The value of the result
lies precisely in its \emph{model independence}.

\section{Scale invariance of the pattern versus occupation}
No scale enters the geometric value $p_0=2/9$ or the ratios: the $C_3$-in-$A_5$ geometry knows
neither mass nor energy, it yields pure numbers -- the angle and the denominator-$9$ weights.
The \emph{pattern layer} -- angle plus ratios -- is therefore a property of the topology, not of
a particular energy range. It holds wherever the $Z_3$/$A_5$ doubling occurs; the lepton sector
is \emph{one} realisation, namely the one in which the scale $E_0$ (Dok.~292) places the absolute
position exactly on the measured $e,\mu,\tau$.

This is to be strictly separated from the \emph{occupation layer}: whether states actually sit at
a given scale. The geometry says \enquote{\emph{if} a $Z_3$ triplet is realised here, then with
this $\theta$ and these ratios} -- it does not say \enquote{here is one}. At another scale the
same circulant pattern stands without occupied states having to lie there. This is the same logic
as the tower $D(4+3n)$ and the emptiness theorem (Dok.~285): the pattern is present everywhere,
the occupation is a separate, scale-bound and dynamics-dependent question; an empty window carries
the same pattern as an occupied one, only without particles in it. From the scale invariance of
the pattern there follows therefore \emph{no} prediction of new particles at other scales -- the
universality concerns the structure, not the occupation.

\section{Status and open edge}
What is secured: The icosahedral doubling structure produces $\theta=2/9$ parameter-free,
forward, exactly and icosahedron-specifically, in independent agreement with the phenomenological
Koide value; the amplitude $\delta=(7-3\varphi)/9$ stems from the same redistribution. Thereby the
long-open question -- whether the $T^4/Z_3$ structure forces $\theta=2/9$ -- has, for the first
time, a concrete, exact, geometric answer within the already established $A_5$ doubling (Dok.~285).

What remains open as a derivation step (and explicitly not as an objection to the convergence):
the explicit mechanism linking the geometric redistribution weights $(2/9,\,(5-3\varphi)/9,\,
(2+3\varphi)/9)$ to the physical mass ratios. This bridge need not run through Koide but can be
led directly through the masses or the mode structure; it is the subject of further work. The
operational role of the $5$-fold in the phase channel is consistent with Dok.~291: $R_5$ sets a
phase $\pi/2$, the value Dok.~291 associates with $\chi$.
